\chapter{Minacce alla validità}
\label{cap:validita}

Nessun esperimento è privo di limiti. Le minacce note sono classificate secondo le quattro categorie usuali e accompagnate dalla contromisura adottata.

\begin{description}
  \item[Interna] un fattore diverso da quello studiato può causare la differenza osservata.
  \item[Di costrutto] la misura può non rappresentare ciò che si vuole misurare.
  \item[Esterna] i risultati possono non valere fuori dall'ambiente di prova.
  \item[Delle conclusioni] le conclusioni statistiche possono essere sbagliate.
\end{description}

% longtable e non tabularx: con le minacce emerse in F3 la tabella non sta più in una pagina.
\newcolumntype{Q}{>{\raggedright\arraybackslash}p{\dimexpr(\linewidth-2.3cm-4\tabcolsep)/2\relax}}
\begin{longtable}{@{}>{\raggedright\arraybackslash}p{2.3cm}QQ@{}}
\caption{Minacce alla validità e contromisure.}
\label{tab:limiti} \\
\toprule
\textbf{Tipo} & \textbf{Minaccia} & \textbf{Contromisura} \\
\midrule
\endfirsthead
\toprule
\textbf{Tipo} & \textbf{Minaccia} & \textbf{Contromisura} \\
\midrule
\endhead
\bottomrule
\endlastfoot
Interna & Sistema e generatore di carico girano sullo stesso computer e si contendono CPU e memoria & Soglia di CPU fissata nella fase F4; run saturi o senza memoria scartati \\
Interna & Un run può lasciare tracce (code non vuote, pod degradati) che influenzano il successivo & Reset verificato prima di ogni run; ordine casuale dei run \\
Interna & Dopo un crash Kubernetes ritarda i riavvii successivi dello stesso container (10, 20, 40~s e così via fino a 5 minuti) e azzera il ritardo solo dopo 10 minuti senza crash: due run con guasto di fila sullo stesso pod avrebbero riavvii diversi & Il reset ricrea i pod prima di ogni run, così ogni guasto colpisce un container senza crash precedenti \\
Interna & Il guasto di un componente può dipendere da dove Kubernetes ha collocato i pod & Regole di piazzamento esplicite (D0006); topologia verificata al reset \\
Di costrutto & Un generatore a modello chiuso smetterebbe di inviare durante il guasto, nascondendolo & Carico a modello aperto; ritmo inviato controllato entro $\pm5\%$ \\
Di costrutto & Un pod ``pronto'' non significa servizio ripristinato & Ripristino definito come 5~s consecutivi disponibili in accettazione e completamento \\
Di costrutto & La soglia del \parametro{soglia_secondo_disponibile_percentuale}\% al secondo è severa a ritmi bassi & Si riportano anche le richieste fallite per guasto, indipendenti dalla soglia \\
Di costrutto & Su AWS gli istanti di ricezione e scrittura vengono da nodi con orologi diversi & Sfasamento verificato e dichiarato prima delle sessioni (D0005) \\
Di costrutto & In locale l'orologio della macchina virtuale di Docker può scostarsi da quello del computer di decine o centinaia di millisecondi, e lo scarto cambia nel tempo & Ogni misura sottrae istanti di un solo orologio: il fermo dell'accettazione usa l'istante del guasto preso sul computer, come il log del client; quello del completamento l'istante letto nel nodo, come il database. Lo scarto tra i due orologi si controlla a ogni guasto \\
Di costrutto & L'istante del guasto non è esatto: il comando che manda il segnale dal nodo impiega da pochi millisecondi a oltre cento su un computer carico & La connessione al nodo si apre prima del guasto; si registrano gli istanti subito prima e subito dopo il segnale, e la loro distanza si dichiara come incertezza del fermo \\
Esterna & In locale i nodi sono container sullo stesso hardware: la perdita di un nodo è simulata & Parte degli esperimenti ripetuta su AWS con macchine reali (RQ6). La perdita di un nodo è provata su AWS solo se si esegue il guasto F6 (facoltativo): altrimenti il limite resta non mitigato \\
Esterna & Il servizio è molto più semplice di uno reale & Le conclusioni riguardano i meccanismi (rilevamento, riparazione, smaltimento), non i valori assoluti \\
Esterna & Il control plane è unico e funge anche da punto d'ingresso: se cadesse, nessun pod verrebbe riavviato e gli ordini non entrerebbero & Non viene mai guastato; gli esperimenti riguardano l'applicazione, non Kubernetes. In produzione il control plane è replicato \\
Delle conclusioni & Un risultato singolo può essere casuale & 10 ripetizioni per i run con guasto; intervalli di confidenza bootstrap \\
Delle conclusioni & Si possono adattare le spiegazioni dopo aver visto i dati & Previsioni registrate prima delle misure (tag \texttt{m1}) \\
\end{longtable}
